% TEMPLATE-MILC-v3.tex - plantilla maestra para todas las guias MILC v3
% Ver scripts/generadores/build-guia-milc.py para la lista completa de
% claves esperadas. El script valida que no queden placeholders sin
% reemplazar antes de escribir el archivo final.

\documentclass[11pt,letterpaper]{article}
\usepackage[margin=1.75cm]{geometry}
\usepackage[table]{xcolor}
\usepackage{fontspec}
\usepackage[spanish]{babel}
\usepackage{tikz}
\usepackage[most]{tcolorbox}
\usepackage{tabularx}
\usepackage{array}
\usepackage{fancyhdr}
\usepackage{titlesec}
\usepackage{ragged2e}
\usepackage{enumitem}
\usepackage{microtype}

\setmainfont{Helvetica Neue}
\setsansfont{Helvetica Neue}
\setlength{\parindent}{0pt}
\setlength{\parskip}{6pt}
\hyphenpenalty=200
\exhyphenpenalty=200
\pretolerance=400
\tolerance=2000
\setlength{\emergencystretch}{1em}
\renewcommand{\arraystretch}{1.25}
\setlist{leftmargin=5mm,itemsep=1mm,topsep=1mm}
\newcolumntype{Y}{>{\RaggedRight\arraybackslash}X}

\definecolor{milcMagenta}{HTML}{E6007E}
\definecolor{milcVino}{HTML}{5A0038}
\definecolor{milcTurquesa}{HTML}{0093A5}
\definecolor{milcVerde}{HTML}{58A923}
\definecolor{milcMostaza}{HTML}{E5B400}
\definecolor{milcOcre}{HTML}{B86600}
\definecolor{milcNegro}{HTML}{241816}
\definecolor{milcGris}{HTML}{F7F7F4}
\definecolor{milcLinea}{HTML}{E8E2DD}
\definecolor{milcGradePrimary}{HTML}{0066FF}
\definecolor{milcGradeDark}{HTML}{003D99}
\definecolor{milcGradeSoft}{HTML}{D6E8FF}
\definecolor{milcGradeAccent}{HTML}{CFE7FF}
\definecolor{milcGradeText}{HTML}{FFFFFF}
\color{milcNegro}

\pagestyle{fancy}
\fancyhf{}
\lhead{\small Tecnología e Informática · Grado 11}
\rhead{\small Guía 2 · MILC v3}
\cfoot{\small\thepage}
\renewcommand{\headrulewidth}{0pt}

\newtcolorbox{titlebox}[1]{
  enhanced,colback=#1,colframe=#1,arc=7mm,boxrule=0pt,
  left=7mm,right=7mm,top=3mm,bottom=3mm,
  before skip=8pt,after skip=8pt,
  fontupper=\sffamily\bfseries\Large\color{white}
}

\newtcolorbox{softbox}[3]{
  enhanced,colback=#2,colframe=#1,borderline west={4pt}{0pt}{#1},
  arc=2mm,boxrule=.4pt,left=4mm,right=4mm,top=3mm,bottom=3mm,
  before skip=6pt,after skip=8pt,title={#3},coltitle=white,
  colbacktitle=#1,fonttitle=\sffamily\bfseries,fontupper=\RaggedRight,
  attach boxed title to top left={xshift=3mm,yshift=-2mm},
  boxed title style={arc=2mm, boxrule=0pt}
}

\newtcolorbox{drawbox}[2]{
  enhanced,colback=white,colframe=#1,arc=4mm,boxrule=1pt,
  left=5mm,right=5mm,top=4mm,bottom=4mm,before skip=8pt,after skip=8pt,
  title={#2},coltitle=white,colbacktitle=#1,fonttitle=\sffamily\bfseries,
  fontupper=\RaggedRight,
  attach boxed title to top left={xshift=4mm,yshift=-2mm},
  boxed title style={arc=3mm, boxrule=0pt}
}

\newtcolorbox{infoband}[1]{
  enhanced,colback=#1!8,colframe=#1!50,arc=2mm,boxrule=.5pt,
  left=4mm,right=4mm,top=2mm,bottom=2mm,before skip=4pt,after skip=4pt,
  fontupper=\small\RaggedRight
}

% Puente narrativo entre secciones (contrato editorial Conéctate).
% Da continuidad: "Acabas de X. Ahora vas a Y porque Z."
% Visualmente: banda izquierda en color del grado, texto en itálica pequeña.
\newtcolorbox{puentebox}{
  enhanced,colback=milcGradeSoft,colframe=milcGradeDark,
  borderline west={3pt}{0pt}{milcGradeDark},
  arc=0mm,boxrule=0pt,left=5mm,right=4mm,top=2.5mm,bottom=2.5mm,
  before skip=4pt,after skip=8pt,
  fontupper=\itshape\small\color{milcNegro}\RaggedRight
}
\newcommand{\puente}[1]{%
  \begin{puentebox}%
    \textcolor{milcGradeDark}{\textbf{\textrightarrow}}\hspace{2mm}#1%
  \end{puentebox}%
}

\newcommand{\linead}[1]{\noindent\color{milcLinea}\rule{#1}{0.5pt}\color{milcNegro}}
\newcommand{\respuesta}[1]{\par\vspace{2mm}\foreach \n in {1,...,#1}{\noindent\color{milcLinea}\rule{\linewidth}{0.5pt}\par\vspace{5mm}}\color{milcNegro}}
\newcommand{\checkbox}{\textcolor{milcGradeDark}{\fbox{\phantom{x}}}\hspace{5pt}}

\newcommand{\phasecard}[4]{
\begin{tcolorbox}[enhanced,colback=#3,colframe=#2,arc=3mm,boxrule=.5pt,height=3.05cm,valign=center,halign=center,left=1.5mm,right=1.5mm,top=2mm,bottom=2mm]
{\sffamily\bfseries\fontsize{22}{24}\selectfont\textcolor{#2}{#1}}\par
{\sffamily\bfseries\small #4}
\end{tcolorbox}
}

\begin{document}
\thispagestyle{empty}
\begin{tikzpicture}[remember picture,overlay]
  \fill[white] (current page.south west) rectangle (current page.north east);
  \path[fill=milcGradePrimary,rounded corners=16mm]
    ([xshift=.55cm,yshift=-.55cm]current page.north west) rectangle
    ([xshift=-.55cm,yshift=3.65cm]current page.south east);

  \node[anchor=north west,text=milcGradeText] at ([xshift=2.0cm,yshift=-1.85cm]current page.north west)
    {\sffamily\bfseries\fontsize{14}{16}\selectfont GRADO};
  \node[anchor=north west,text=milcGradeText,opacity=.82] at ([xshift=2.0cm,yshift=-2.75cm]current page.north west)
    {\sffamily\bfseries\fontsize{9}{11}\selectfont SERIE GUÍAS · TECNOLOGÍA E INFORMÁTICA};

  \begin{scope}[shift={([xshift=-3.05cm,yshift=-2.55cm]current page.north east)}]
    \fill[white,opacity=.80] (-.62,-.52) rectangle (.62,.52);
    \draw[milcGradeText,line width=1pt] (-.62,-.52) rectangle (.62,.52);
    \fill[milcGradeDark] (-.42,-.38) rectangle (-.22,.06);
    \fill[milcGradeAccent] (-.10,-.38) rectangle (.10,.24);
    \fill[milcGradePrimary!60!black] (.22,-.38) rectangle (.42,.38);
  \end{scope}

  \node[anchor=west,text=milcGradeText] at ([xshift=1.9cm,yshift=-12.85cm]current page.north west)
    {\sffamily\bfseries\fontsize{124}{124}\selectfont 11};
  \draw[milcGradeText,line width=1.7pt]
    ([xshift=4.52cm,yshift=-11.68cm]current page.north west) circle (.13cm);

  \node[anchor=west,text=milcGradeText,text width=8.7cm,align=left] at ([xshift=10.35cm,yshift=-11.6cm]current page.north west)
    {
     {\sffamily\bfseries\fontsize{19}{22}\selectfont Identidad visual ---\\del grafismo ancestral\\a la paleta propia}\\[4mm]
     {\sffamily\bfseries\fontsize{23}{26}\selectfont Guía 2-11-TIC}\\[2mm]
     {\sffamily\fontsize{13}{16}\selectfont Período 1 · Presencia y marca digital}\\[5mm]
     {\sffamily\bfseries\fontsize{11}{13}\selectfont Institución Educativa Sor María Juliana}\\[1mm]
     {\sffamily\fontsize{10}{12}\selectfont Autor: Dr. Álvaro Cárdenas Orozco}
    };

  \path[fill=milcGradeSoft,rounded corners=7mm]
    ([xshift=1.35cm,yshift=.85cm]current page.south west) rectangle
    ([xshift=-1.35cm,yshift=4.25cm]current page.south east);
  \draw[milcGradeDark,line width=.65pt,rounded corners=7mm]
    ([xshift=1.35cm,yshift=.85cm]current page.south west) rectangle
    ([xshift=-1.35cm,yshift=4.25cm]current page.south east);
  \node[anchor=center,text=milcNegro,text width=17.0cm,align=center] at ([yshift=2.55cm]current page.south)
    {\sffamily\fontsize{10}{12}\selectfont Metodología MILC · Apertura ancestral · Pensamiento computacional · Triángulo Dussel-Estoicismo-Floridi\\
    \textcolor{milcGradeDark}{\bfseries Producto:} Sistema visual mínimo: 1 paleta (3-5 colores) · 1 par tipográfico · 1 lockup};
\end{tikzpicture}
\mbox{}
\newpage

\section*{Apertura: Identidad visual --- del grafismo ancestral a la paleta propia}

\begin{softbox}{milcGradeDark}{milcGradeSoft}{Saber ancestral}
Mucho antes de los ``manuales de marca'', los pueblos ya tenían identidad visual. La \textbf{orfebrería Quimbaya}, los \textbf{tejidos Wayuu}, las \textbf{molas guna} y los grafismos del Pacífico colombiano transmitían, en una sola pieza, quién era el pueblo que la había hecho. Sin firma, sin texto. Solo geometría, color y proporción. Esa gramática visual sostenía el reconocimiento entre generaciones.
\end{softbox}

\begin{softbox}{milcTurquesa}{milcGris}{Saber contemporáneo}
Las marcas modernas usan exactamente el mismo principio en formato digital: una \textbf{paleta} de 3 a 5 colores (no más), un \textbf{par tipográfico} (una fuente para titulares y una para cuerpo) y un \textbf{lockup} (cómo se firma el nombre y el logo). Esa pieza mínima identifica tu marca antes de que el espectador lea una palabra.
\end{softbox}

\begin{softbox}{milcMostaza}{milcGris}{Pregunta-puente}
\textit{¿Qué señales visuales hacen que reconozcas inmediatamente a una familia, a un equipo del barrio o a tu colegio sin leer su nombre? ¿Por qué esas señales funcionan?}
\end{softbox}

\begin{softbox}{milcVerde}{milcGris}{Saber y saber hacer}
Al terminar podrás: (1) \textbf{identificar} gramáticas visuales en tu entorno cercano (ancestrales, populares, comerciales); (2) \textbf{analizar} el sistema visual de tres marcas que admiras; (3) \textbf{crear} tu propio sistema visual mínimo coherente con tu manifiesto; (4) \textbf{evaluar} el sistema de un compañero con criterios de coherencia y diferenciación.
\end{softbox}

\small
\begin{tabularx}{\textwidth}{>{\bfseries}p{3.0cm}Y}
\rowcolor{milcGradePrimary!12}Autor & Dr. Álvaro Cárdenas Orozco\\
DBA & Comunicación profesional en entornos digitales (MEN, Lineamientos T\&I)\\
Referentes & Floridi (infoética) · Dussel (estética de la liberación) · Estoicismo\\
Producto final & Sistema visual mínimo: 1 paleta (3-5 colores) · 1 par tipográfico · 1 lockup\\
\end{tabularx}
\normalsize

\puente{Antes de empezar, mira el viaje completo. Hoy le pondrás \emph{cara} a la marca que definiste la semana pasada --- la promesa que escribiste necesita una piel visual que la sostenga.}

\begin{titlebox}{milcGradeDark}
Ruta MILC: así vamos a aprender
\end{titlebox}
\begin{tabularx}{\textwidth}{>{\centering\arraybackslash}X>{\centering\arraybackslash}X>{\centering\arraybackslash}X>{\centering\arraybackslash}X}
\phasecard{1}{milcVerde}{milcGris}{Escuta\\[-1mm]\small Observo, escucho y pregunto.} &
\phasecard{2}{milcTurquesa}{milcGris}{Sistematización\\[-1mm]\small Pensamiento computacional.} &
\phasecard{3}{milcMagenta}{milcGris}{Praxis\\[-1mm]\small Construyo y aplico.} &
\phasecard{4}{milcVino}{milcGris}{Evaluación\\[-1mm]\small Reflexión liberadora.}
\end{tabularx}


\puente{Toda identidad visual nace de \emph{ver primero}. Antes de diseñar la tuya, vas a entrenar el ojo: cómo se reconocen pueblos, oficios y barrios sin necesidad de palabras. Sin ese entrenamiento, cualquier paleta que elijas será aleatoria.}

\begin{titlebox}{milcVerde}
Fase 1 · Escuta
\end{titlebox}

\begin{softbox}{milcVerde}{milcGris}{Escena para escuchar}
\textbf{Actividad 1 · IDENTIFICA --- Gramáticas visuales en tu entorno} (15 min · individual). Sin moverte de donde estás, observa con atención \textbf{tres ejemplos} de identidad visual cercana: (1) un objeto de tradición ancestral o popular (tejido, cerámica, ropa de fiesta), (2) un negocio del barrio (panadería, peluquería, taller), (3) una marca global presente en tu cotidianidad (zapatillas, app, café). Para cada uno: ¿qué colores usa? ¿qué formas? ¿qué transmite sin palabras?
\end{softbox}

\checkbox Encontré un ejemplo de identidad visual ancestral o popular en mi entorno.\\
\checkbox Encontré un ejemplo de identidad visual de un negocio local que reconozco.\\
\checkbox Encontré un ejemplo de identidad visual global presente en mi cotidianidad.

\begin{infoband}{milcVerde}
\textbf{Lo que va al cuaderno (Actividad 1 · IDENTIFICA).} Título: «Actividad 1 --- Gramáticas visuales en mi entorno». \textbf{Formato:} tabla de 4 columnas (Ejemplo | Colores principales | Formas y patrones | Qué transmite). \textbf{Extensión:} 3 filas, una por ejemplo. \textbf{Sabes que terminaste cuando} cada fila menciona colores específicos (no ``varios''), formas concretas (no ``diseños bonitos'') y un mensaje que el sistema visual transmite sin texto.
\end{infoband}

\puente{Ya viste cómo otros transmiten identidad sin palabras. Ahora vamos a entender los \textbf{tres componentes} que toda marca digital usa para hacerlo: paleta, tipografía y lockup. Tres piezas, una decisión coherente.}

\begin{titlebox}{milcTurquesa}
Fase 2 · Sistematización con pensamiento computacional
\end{titlebox}

\begin{softbox}{milcTurquesa}{milcGris}{Cuatro pilares aplicados al tema}
Todo sistema visual moderno se compone de tres piezas mínimas. Vamos a descomponerlas con los pilares del pensamiento computacional para que puedas diseñar el tuyo con criterio.
\end{softbox}

\small
\begin{tabularx}{\textwidth}{>{\bfseries\RaggedRight\arraybackslash}p{3.6cm}Y}
\rowcolor{milcTurquesa!90}\color{white}Pilar & \color{white}\bfseries Aplicación al tema\\
1. Descomposición & Una identidad visual se descompone en: (1) \textbf{paleta} (3-5 colores con función clara), (2) \textbf{tipografía} (1 fuente display para titulares + 1 fuente cuerpo para texto), (3) \textbf{lockup} (la firma del nombre con su tratamiento visual).\\
2. Reconocimiento de patrones & Toda identidad fuerte sigue patrones de contraste y proporción: un color dominante (\textasciitilde 60\%), uno de soporte (\textasciitilde 30\%), uno de acento (\textasciitilde 10\%); un tipo display contrastado contra el cuerpo; espacio en blanco generoso. Cambia las proporciones y se rompe.\\
3. Abstracción & Lo esencial de un sistema visual cabe en \emph{una pieza}: si tu lockup, en blanco y negro y del tamaño de una uña, todavía dice quién eres, has abstraído bien.\\
4. Algoritmo conceptual & Pasos: (1) elige una emoción principal de tu marca (calidez, rigor, frescura, etc.); (2) escoge la paleta que la transmita; (3) escoge el par tipográfico; (4) construye el lockup que firma todo.\\
\end{tabularx}
\normalsize

\begin{drawbox}{milcTurquesa}{Anatomía de un sistema visual mínimo}
\textbf{Paleta:} 3-5 colores con jerarquía (uno dominante, uno secundario, 1-3 acentos). Evita colores neón puros para texto. \\[1mm] \textbf{Tipografía:} 1 display + 1 cuerpo. Si tu display es fuerte, tu cuerpo debe ser neutro --- nunca dos display al mismo tiempo. \\[1mm] \textbf{Lockup:} la firma de tu nombre con tratamiento fijo (color, fuente, espacios). Debe funcionar a 16 píxeles y a 2 metros de pared.
\end{drawbox}

\begin{softbox}{milcMostaza}{milcGris}{Errores comunes y su corrección}
(1) Más de 5 colores: la identidad se diluye. (2) Dos fuentes display compitiendo: se siente caótico. (3) Color neón saturado en texto: ilegible. (4) Copiar la paleta exacta de una marca global: tu marca pierde identidad.
\end{softbox}

\begin{infoband}{milcTurquesa}
\textbf{Lo que va al cuaderno (Actividad 2 · ANALIZA).} Título: «Actividad 2 --- Tres sistemas visuales que admiro». \textbf{Formato:} 3 fichas, una por marca, con 4 datos cada una: Nombre · Paleta (3-5 códigos HEX) · Par tipográfico · Qué transmite (1 frase). \textbf{Sabes que terminaste cuando} las 3 marcas son distintas entre sí (no las 3 deportivas), y la frase ``qué transmite'' es específica (no ``se ve cool'').
\end{infoband}

\puente{Tienes el vocabulario y los referentes. Toca decidir: ¿qué paleta, qué tipografía y qué lockup hablan por la marca que escribiste la semana pasada? Esto es lo que vas a entregar hoy.}

\begin{titlebox}{milcMagenta}
Fase 3 · Praxis con pensamiento computacional
\end{titlebox}

\begin{softbox}{milcMagenta}{milcGris}{Construyendo el producto con los cuatro pilares}
\textbf{Actividad 3 · CREA --- Tu sistema visual mínimo} (30 min · individual). Tienes el vocabulario y los referentes. Hora de tomar decisiones para tu marca. Recuerda: tu manifiesto (Guía 1) prometió algo --- tu sistema visual debe sostener esa promesa.
\end{softbox}

\small
\begin{tabularx}{\textwidth}{>{\bfseries\RaggedRight\arraybackslash}p{3.6cm}Y}
\rowcolor{milcMagenta!90}\color{white}Pilar & \color{white}\bfseries Aplicación al producto\\
Descomposición & Tu sistema se descompone en 3 piezas: paleta + tipografía + lockup. Cada una contesta una pregunta: ¿qué emoción transmite? ¿qué jerarquía organiza el texto? ¿cómo firma todo?\\
Patrones & Sigue el patrón 60-30-10 en color: un dominante, un secundario, un acento. Sigue el patrón display + cuerpo en tipografía. Sigue el patrón ``nombre + tratamiento'' en lockup.\\
Abstracción & Cada decisión expresa una sola idea. Si no puedes decir en una frase corta por qué elegiste ese color, ese tipo o ese lockup, vuelve atrás.\\
Algoritmo de armado & Algoritmo: (1) elige la emoción principal; (2) prueba 3 paletas en coolors.co; (3) prueba 3 pares tipográficos en fonts.google.com; (4) dibuja a mano 3 versiones de tu lockup. Elige la combinación más coherente.\\
\end{tabularx}
\normalsize

\begin{drawbox}{milcMagenta}{Checklist del sistema visual antes de cerrar el cuaderno}
\checkbox Paleta de 3 a 5 colores con códigos HEX anotados. \\[1mm] \checkbox Par tipográfico nombrado (ej.: Bricolage Grotesque + Inter). \\[1mm] \checkbox Lockup dibujado a mano o en herramienta digital. \\[1mm] \checkbox Justificación de 1 línea por cada decisión (color, tipo, lockup). \\[1mm] \checkbox Tu manifiesto de la Guía 1 se refleja coherente en este sistema.
\end{drawbox}

\begin{softbox}{milcMagenta}{milcGris}{Plantilla de trabajo}
\textbf{Mi paleta.} \textit{HEX codes + para qué sirve cada color (dominante/secundario/acento).} \\[1mm] \textbf{Mi par tipográfico.} \textit{Display + cuerpo, nombres y dónde se encuentran.} \\[1mm] \textbf{Mi lockup.} \textit{Cómo firma tu marca: dibujo, tratamiento, tamaños.} \\[1mm] \textbf{Por qué estas decisiones.} \textit{1 frase por elemento.}
\end{softbox}

\begin{infoband}{milcMagenta}
\textbf{Lo que va al cuaderno (Actividad 3 · CREA).} Título: «Actividad 3 --- Mi sistema visual · versión 1». \textbf{Formato:} una página con 4 secciones (Paleta · Tipografía · Lockup · Justificación), cada una con datos concretos y una frase de decisión. \textbf{Extensión:} una página de cuaderno. \textbf{Sabes que terminaste cuando} los 4 elementos están presentes, cada decisión tiene su frase de justificación, y tu lockup se puede dibujar a mano sin guías.
\end{infoband}

\puente{Lo que sigue es el resumen de \emph{qué} entregas y \emph{cómo se sabe} si tu sistema visual es coherente. El primer criterio: que un extraño pueda mirarlo y sentir tu marca sin que tú la expliques.}

\begin{titlebox}{milcMostaza}
Producto de la sesión
\end{titlebox}
\textbf{Sistema visual mínimo en una página: paleta + par tipográfico + lockup}.
\begin{softbox}{milcMostaza}{milcGris}{Criterios de calidad}
(1) La paleta tiene 3-5 colores con jerarquía clara. (2) El par tipográfico es uno display + uno cuerpo, no dos display. (3) El lockup tiene tratamiento fijo y funciona a tamaño chico. (4) Cada decisión tiene una frase de justificación que la conecta con tu manifiesto. (5) Un extraño puede mirar el sistema y sentir tu marca sin que tú la expliques.
\end{softbox}

\puente{Antes de cerrar, mira tu sistema visual desde las cinco dimensiones humanas. Las decisiones estéticas tienen consecuencias personales, ciudadanas, locales --- aunque a primera vista no lo parezca.}

\begin{titlebox}{milcVino}
Fase 4 · Evaluación liberadora — cinco dimensiones
\end{titlebox}

\begin{softbox}{milcVino}{milcGris}{Sentido de la evaluación}
Evaluar no es solo revisar una nota. Es reconocer qué aprendí, qué cambió en mí, qué emociones manejé, cómo me relaciono con la ciudad y la comunidad, y qué le dejaré a quien viene detrás.
\end{softbox}

\textbf{Mi reflexión liberadora en cinco dimensiones.} Completa cada frase iniciada:

\small
\begin{tabularx}{\textwidth}{>{\bfseries\RaggedRight\arraybackslash}p{3.4cm}Y>{\RaggedRight\arraybackslash}p{4.6cm}}
\rowcolor{milcVino}\color{white}Dimensión & \color{white}\bfseries Mi respuesta (frame starter) & \color{white}\bfseries Algo que descubrí\\
1. Personal & Lo que más cambió en mí fue \linead{6.5cm} & \linead{4.2cm}\\
2. Emocional & La emoción más fuerte fue \linead{2.5cm} y la regulé \linead{3cm} & \linead{4.2cm}\\
3. Ciudadana & En democracia esto importa porque \linead{6cm} & \linead{4.2cm}\\
4. Local (Cartago) & En mi barrio sería útil para \linead{6.5cm} & \linead{4.2cm}\\
5. Intergeneracional & A mi abuela/abuelo le diría \linead{6.5cm} & \linead{4.2cm}\\
\end{tabularx}
\normalsize

\begin{infoband}{milcVino}
\textbf{Para la próxima sustentación.} Vuelve a esta tabla en seis meses. Verás que las respuestas cambian — no porque lo aprendido cambie, sino porque tú cambias. La evaluación liberadora es una conversación contigo a través del tiempo.
\end{infoband}


\puente{Y para terminar, tres voces sobre la imagen. Dussel sobre la belleza descolonizada, Marco Aurelio sobre la belleza propia, Floridi sobre la responsabilidad visual en la infoesfera. Cada uno te ofrece un lente distinto para mirar el sistema que acabas de crear.}

\begin{titlebox}{milcGradeDark}
Triángulo de pensamiento · cierre filosófico
\end{titlebox}

\begin{softbox}{milcVino}{milcGris}{Dussel · lente del nosotros}
\textit{``La belleza estética nace donde una voz silenciada encuentra forma.''}

Tu sistema visual puede repetir las paletas de moda global (azules tech, naranjas startup) o puede atreverse a recuperar colores y formas de tu lugar --- los amarillos del Valle, los rojos del Pacífico, los grafismos del oficio. La descolonización estética empieza por mirar antes de copiar.

\textbf{¿Mi sistema visual incluye señales de mi lugar y mi gente, o copia el estándar global? ¿A qué voz le doy forma cuando elijo estos colores?}
\end{softbox}
\linead{\linewidth}\\[1mm]
\linead{\linewidth}

\begin{softbox}{milcMostaza}{milcGris}{Estoicismo · Marco Aurelio · cuidado interior}
\textit{``Toda cosa tiene su belleza propia, pero no todos la ven.''}

La tentación más fuerte en diseño visual es seguir tendencias. Pero las tendencias pasan; la belleza propia de tu marca se descubre cuando dejas de mirar lo que está de moda y miras lo que de verdad transmite tu promesa. Eso depende de ti.

\textbf{¿Estoy eligiendo lo que de verdad expresa mi marca o estoy persiguiendo una estética que vi en otra parte? ¿Cuál es mi belleza propia?}
\end{softbox}
\linead{\linewidth}\\[1mm]
\linead{\linewidth}

\begin{softbox}{milcTurquesa}{milcGris}{Floridi · lente de la infoesfera}
\textit{``Cada acto de diseño es un acto informacional: añade o resta a la infoesfera visual común.''}

Tu sistema visual va a aparecer en feeds, presentaciones, perfiles. Va a competir por atención en una infoesfera ya saturada. La pregunta no es ``¿llama la atención?'' sino ``¿añade algo útil o agrega ruido visual?''. El diseño tiene ética informacional.

\textbf{¿Mi sistema visual aporta claridad a quien lo mira o le suma ruido a un mundo ya saturado? ¿Qué responsabilidad tengo sobre la atención de quienes lo verán?}
\end{softbox}
\linead{\linewidth}\\[1mm]
\linead{\linewidth}

\begin{titlebox}{milcVino}
Compromiso final
\end{titlebox}

\small
\begin{tabularx}{\textwidth}{>{\RaggedRight\arraybackslash}p{3.0cm}YYY}
\rowcolor{milcVino}\color{white}\bfseries Criterio & \color{white}\bfseries Lo logré & \color{white}\bfseries En proceso & \color{white}\bfseries Necesito apoyo\\
Comprensión del tema & Explico ideas centrales con mis palabras. & Reconozco ideas pero falta precisión. & Necesito repasar conceptos base.\\
Pensamiento computacional & Apliqué los 4 pilares al diseñar mi producto. & Apliqué algunos pilares. & Necesito releer la fase 2.\\
Aplicación y producto & Mi producto cumple los criterios y se puede revisar. & Producto funcional con ajustes pendientes. & Aún en construcción.\\
Reflexión 5 dimensiones & Respondí cada dimensión con honestidad. & Respondí algunas con detalle. & Mis respuestas son muy generales.\\
Triángulo de pensamiento & Conecté las 3 voces con mi trabajo real. & Conecté al menos una. & No relaciono las voces con mi proceso.\\
\end{tabularx}
\normalsize

\textbf{Hoy entendí que:}\\[1mm]
\linead{\linewidth}\\[1mm]
\linead{\linewidth}

\textbf{Mi compromiso para la próxima sesión es:}\\[1mm]
\linead{\linewidth}\\[1mm]
\linead{\linewidth}

\begin{softbox}{milcMostaza}{milcGris}{Cita de despedida · Marco Aurelio}
\textit{``El obstáculo en el camino se vuelve el camino. Lo que estorba al camino se vuelve camino.''}\\[1mm]
Cada error de hoy, bien observado, se convierte en pieza del camino que pisarás mañana.
\end{softbox}

\small
\begin{tabularx}{\textwidth}{>{\bfseries}p{3.6cm}Y}
\rowcolor{milcGradeDark!12}Nombre completo: & \linead{10cm}\\
Fecha: & \linead{4cm} \quad Curso/grupo: \linead{4cm}\\
Mi firma: & \linead{9cm}\\
\end{tabularx}
\normalsize

\begin{infoband}{milcGradeDark}
\textbf{Para la próxima semana.} Aplica tu paleta a 3 productos cotidianos: la portada de tu próximo trabajo, tu foto de perfil con un filtro de color coherente, y la firma de tu correo institucional. Ese ensayo te mostrará si tu sistema funciona o necesita ajustes. El próximo lunes traes esos 3 productos al aula.
\end{infoband}

\end{document}
